% Shared section of the technical documents of the six fixed-B modes: the adjoint of the embedded solve.
% \input by each mode document (after \adjointcolumntrue or \adjointpoissontrue, see ../preamble.tex).
\section{Adjoint (embedded API)}\label{sec:adjointapi}
The cotangents of a solve's outputs pulled back to its inputs ($\ln T$ and the cell-centred $\B$), at a cost independent of
how many inputs a caller varies. Design and decisions: \code{design/adjoint/}; results: \code{HSEMHS/reports/adjoint\_phase1-3.md},
\code{adjoint\_phase4-6.md}. Code: \code{src/Adjoint.jl} (API), \code{src/ColumnHSEAdjoint.jl}, \code{src/PoissonTauAdjoint.jl},
\code{src/PoissonZAdjoint.jl}, \code{src/Tau.jl}, \code{src/ViewTau.jl}; the table derivatives are WittBifrostEOS's \code{state\_d}.
{\footnotesize\begin{verbatim}
res = solve_embedded!(ws, lnT, Bx, By, Bz; record = true)  # keeps the tape; any later solve ends it
bar = EquilibriumCotangent(ws)                             # zeros on this rank's slab, (vertical, y, x)
bar.rho .= ...; bar.z .= ...       # any of Pgas, lnP, rho, ne, chi500, tau500, logtau500, z
adj = adjoint_embedded!(ws, bar)   # collective; the tape is kept, so several cotangents work
adj.lnT, adj.Bx, adj.By, adj.Bz    # d(loss)/d ln T and d(loss)/dB on the slab
adj.status                         # :exact, :stationary or :not_stationary; details in adj.info
\end{verbatim}}
\begin{itemize}
\item The call must follow a valid \code{solve\_embedded!} with \code{record = true} on the same workspace; a workspace without
a tape, after a later solve, or whose last solve had a given $n_e$ (\code{lnne}) is refused (\code{InvalidInput}, on every rank).
\code{Pgas} and \code{lnP} cotangents may both be given; they are combined. \code{z\_tau1} is a diagnostic and has no cotangent.
\item Keywords: \code{cg\_tol} ($10^{-10}$; the Poisson adjoint solves), \code{cg\_maxit} (1000), \code{gauss\_newton} (false; true
uses the Gauss--Newton Hessian instead, for comparison), \code{stationary\_tol} ($10^{-6}$).
\ifadjointcolumn
\item \textbf{Column modes:} reverse sweeps of the recorded integrator (\code{RK4Tape} on the height grid, \code{HeunTape} on the
$\tau$ and mass axes: substep counts, the bound that set each Heun step and how often it was quartered, table cells), the top
closure by the implicit function theorem, the alignment with its median reference (\code{gquantile\_vjp}), the gauge, the
$\tau$ integration (\code{tau\_column\_vjp!}, \code{inclined\_tau\_vjp!} with the transposed shear) and the state at the nodes.
\code{adj.status = :exact}. No communication except the reductions of the alignment and the gauge.
\fi
\ifadjointpoisson
\item \textbf{Poisson modes:} the implicit-function derivative at the returned state: one CG solve with the exact Hessian of $\Phi$
(Hessian--vector products by central differences of the exact gradient; relative step $10^{-5}$ on the $\tau$ grid,
\code{HVP\_STEP}), the transposed mixed derivative in $(\ln T,\B)$, the weights through their layer-median columns, the
continuation or fixed part above the cut, the gauge, and the column-HSE start. Directions in which $\Phi$ does not change are
projected out of the solve and their cotangent passed to the start (height grid: the anchor-cell mean of $\ln P$, a bordered
constraint; $\tau$ and mass axes: the uniform height shift and the Nyquist modes of the centred differences). The solve has its
own block-Jacobi preconditioner (the exact Gauss--Newton column blocks with no damping on $\ln P$ and the measured curvature
on $c$); the solver's preconditioners, built for its damped operator, made CG break down. \code{adj.status} is
\code{:stationary} when one undamped Newton step from the returned state predicts a relative decrease of $\Phi$ at most
\code{stationary\_tol} (\code{adj.info["newton\_decrement"]}, solved to CG tolerance $10^{-3}$), else \code{:not\_stationary}.
\item \textbf{\code{polish} = $n$} (\code{PoissonZOptions}, \code{PoissonTauOptions}, default 0; embedding API only, not a
\code{[poisson]} key of the file driver): up to $n$ Newton
steps with the exact Hessian after the Gauss--Newton solve, so that the returned state is the stationary point the adjoint
differentiates. Without it the Gauss--Newton solves can end far from stationarity although their damped decrement is small: on
the $11\times12$ height-grid fixture 80, 300 and 1000 steps gave finite differences 4.4, 7.8 and 10.0 against the adjoint's
10.2, and the anchor mean drifted by $5\times10^{-3}$; 30 steps and 4 Newton steps reach a decrement of $10^{-25}$.
\fi
\end{itemize}

\textbf{Accuracy} (dot-product check against central differences of \code{solve\_embedded!}, Richardson, best over $h$;
\code{test/test\_adjoint.jl}, \code{test/mpi/mpi\_adjoint.jl} on 2 and 3 ranks, which also equal rank 0 alone to $10^{-10}$):
\begin{center}\small
\begin{tabular}{lcc}
\toprule
case & 1 rank & 3 ranks \\ \midrule
\ifadjointcolumn
column-hse-ltau (ptot, $\B$); align none, $\B=0$; $\mu$ 0.8 inclined & $3.8\times10^{-10}$; $1.8\times10^{-9}$; $9.4\times10^{-9}$ & the same \\
column-hse-logm (ptot, $\B$); $\mu$ 0.8 inclined & $9.3\times10^{-10}$; $2.6\times10^{-10}$ & $9.4\times10^{-10}$; $2.5\times10^{-10}$ \\
column-hse-z anchor bottom / top / height; $\mu$ 0.8 inclined & $3.0\times10^{-8}$ / $5.7\times10^{-10}$ / $4.9\times10^{-9}$; $2.3\times10^{-9}$ & the same \\
\fi
\ifadjointpoisson
mhs-poisson-z (30 Gauss--Newton steps + polish) & $6.4\times10^{-8}$ & $3.5\times10^{-8}$ \\
mhs-poisson-ltau, default solver + polish & $3.0\times10^{-8}$ & $1.4\times10^{-9}$ \\
mhs-poisson-ltau $\mu$ 0.8 inclined; block preconditioner & $5.8\times10^{-9}$; $2.3\times10^{-9}$ & $5.8\times10^{-9}$; $2.0\times10^{-9}$ \\
mhs-poisson-logm, default solver + polish & $3.5\times10^{-9}$ & $3.9\times10^{-9}$ \\
\fi
\bottomrule
\end{tabular}
\end{center}
\textbf{Cost} (252$^2$ columns, one rank $\times$ 64 threads, hyades5, \code{test/scripts/adjoint\_benchmark.jl}; the column
times vary by about 30\,\% between runs):
\begin{center}\small
\begin{tabular}{lccc}
\toprule
mode & solve & recorded solve & adjoint \\ \midrule
\ifadjointcolumn
column-hse-ltau (96 levels) & 0.37\,s & 1.24\,s & 0.60\,s \\
column-hse-logm (96 levels) & 0.36\,s & 0.79\,s & 0.63\,s \\
column-hse-z (177 levels) & 0.81\,s & 1.74\,s & 2.51\,s \\
\fi
\ifadjointpoisson
mhs-poisson-ltau (95 levels, polish 3), 1 rank $\times$ 64 threads & 631\,s & 649\,s & 555\,s \\
the same, 8 ranks $\times$ 8 threads on one Zen\,3 node (Sect.~\ref{sec:parallel}) & 67\,s & -- & 52\,s \\
mhs-poisson-z (177 levels, polish 3), 1 $\times$ 64 / 8 $\times$ 8 & 281\,s / 12.8\,s & 275\,s / -- & 9.5\,s / 2.2\,s \\
\fi
\bottomrule
\end{tabular}
\end{center}
Limits: the derivatives at the table's slope breaks (node pairs $10^{-9}$ apart in $\ln T$) are those of the cell the point
falls in, so a finite difference straddling a break does not match; the median weights are exact away from ties.
